% Apresentação concisa dos pontos relevantes do trabalho. Deve ser informativo, apresentando finalidades, metodologia, resultados e conclusões; composto de uma sequência de frases concisas, afirmativas e não de enumeração de tópicos. Deve-se usar o verbo na voz ativa e na terceira pessoa do singular, contendo de 150 a 500 palavras. Deve-se evitar símbolos que não sejam de uso corrente e fórmulas, equações, diagramas etc. que não sejam absolutamente necessários. Após o texto do resumo, recomenda-se que sejam inseridas de 3 a 5 palavras-chave. Xxxxxxx xxxxxxxx xxxxxx xxxxxx xxxxxxxx xxxxxxxxxxxxxxxxxxx. Xxxxxxxxxxx xxxxxxxx xxxxxxx xxxxxxx xxxxxxxx xxxxxxxx xxxxxx xxxxxxxx xxxxxxx xxxxxxxxxxxxxx  xxxxxxxx xxxxxxxxxxxxx xxxxxxxxxxx  xxx xxxxxxxxxxxxxx  xxxxxxxx xxxxxxxxxxxxx xxxxxxxxxxx  xxx xxxxxxxxxxxxxx  xxxxxxxx xxxxxxxxxxxxx xxxxxxxxxxx  xxx xxxxxxxx xxxxxx xxxxxx xxxxxxxx xxxxxxxxxxxxxxxxxxx. Xxxxxxxxxxx xxxxxxxx xxxxxxx xxxxxxx xxxxxxxx xxxxxxxx xxxxxx xxxxxxxx xxxxxx xxxxxx xxxxxxxx xxxxxxxxxxxxxxxxxxx. Xxxxxxxxxxx xxxxxxxx xxxxxxx xxxxxxx xxxxxxxx xxxxxxxx xxxxxx. Xxxxxxxx xxxxxx xxxxxx xxxxxxxx xxxxxxxxxxxxxxxxxxx. Xxxxxxxxxxx xxxxxxxx xxxxxxx xxxxxxx xxxxxxxx xxxxxxxx xxxxxx. Xxxxxxxx xxxxxx xxxxxx xxxxxxxx xxxxxxxxxxxxxxxxxxx. Xxxxxxxxxxx xxxxxxxx xxxxxxx xxxxxxx xxxxxxxx xxxxxxxx xxxxxx.
Este trabalho apresenta o desenvolvimento de uma plataforma robótica móvel a partir da adaptação de um robô aspirador de chão comercial para funcionar como um robô seguidor de linha utilizando controle Proporcional-Integral-Derivativo (PID). Com isto, visa-se criar uma plataforma robótica aberta e documentada, que pode ser replicada e adaptada por outros pesquisadores, estudantes e professores para diversas aplicações educacionais e práticas. A plataforma é baseada em um Raspberry Pi com sistema Linux Ubuntu e Robot Operating System (ROS2), responsável pelo processamento e controle, e um microcontrolador ESP32 dedicado à leitura dos sensores de linha. O robô aspirador foi modificado para permitir integração com os componentes eletrônicos e software, tornando-se uma ferramenta para a aplicação em projetos de robótica aplicada nos mais diversos setores. A proposta inclui a implementação da comunicação entre os dispositivos, proporcionando um ambiente prático e acessível para experimentos e aprendizado. Por fim, será elaborada documentação para operação e modificação da plataforma.

% Separe as palavras-chave por ponto
\palavraschave{Robótica, seguidor de linha, PID, Raspberry Pi, ROS2, ESP32}
